\documentclass[11pt]{article}
\usepackage{fontspec}
\usepackage[margin=1in]{geometry}
\usepackage{enumitem}
\usepackage{hyperref}

\setmainfont{Latin Modern Roman}
\setlist[itemize]{leftmargin=1.4em}

\title{Prompt 07 Portal Matters Refactor}
\author{Justicia Administrativa Live Site}
\date{2026-05-13}

\begin{document}
\maketitle

\section*{Scope Lock}
Prompt 07 was limited to \texttt{001-apps/001-site/**}. The live
\texttt{new.justiciadministrativa.com} website continues to be served from
\texttt{001-apps/001-site}. \texttt{001-apps/002-web} was not modified and
remains prototype/reference material only. ArgumentaLab, Laravel, Filament,
\texttt{002-clients/**}, mobile, generated exports, root docs, and root files
were not modified.

\section*{Prompt 06 Handoff}
Prompt 06 commit
\texttt{0edf1075bcf19fd1ba5f31aa806ed69a0d79ef76} was present and contained
only \texttt{001-apps/001-site/**} paths. It split the portal messages view and
left \texttt{matters\_view.dart} as the longest remaining near-threshold
live-site Dart file.

\section*{Matters Ownership Audit}
\texttt{matters\_view.dart} was selected because it was dirty only with
coherent, matter-specific, site-local changes. No existing
\texttt{portal/matters/**} files were present, so the extraction could proceed
without absorbing unrelated site-local work.

\section*{Component Boundaries}
\begin{itemize}
  \item \texttt{matters\_view.dart} remains the stateful page boundary. It owns
  matter loading, fuzzy-search indexing, search debounce, status state, and PDF
  preview orchestration.
  \item \texttt{matters\_page\_header.dart} owns the page title and desktop
  primary action.
  \item \texttt{matters\_filter\_bar.dart} owns the responsive search field and
  status chips.
  \item \texttt{matters\_list.dart} owns list filtering, loading skeletons,
  empty state, and sliver list composition.
  \item \texttt{matter\_card.dart} owns matter summary presentation, status
  facts, context-menu actions, report-preview callback wiring, and animation.
\end{itemize}

\section*{UI and Accessibility}
The refactor preserved the matter list hierarchy, pinned filter surface, fuzzy
search behavior, status filters, report preview affordance, context menu,
card animation, desktop/mobile layout, and brand tokens. A semantic label was
added to each matter card so screen-reader users get title, expediente, and
status context without changing the visual design.

\section*{Architecture Notes}
This slice uses composition and presentational component extraction rather than
new domain abstractions. \texttt{Matter} remains the existing UI data contract.
No new ADTs, repositories, services, or global view models were introduced.
Actor-model abstractions were not needed: asynchronous work is limited to the
page boundary loading matters and opening report previews through explicit
callbacks.

\section*{Validation Results}
\begin{itemize}
  \item \texttt{dart format} on touched Dart files: passed after an escalated
  run for the shared Flutter cache.
  \item \texttt{flutter analyze}: passed with no issues.
  \item \texttt{flutter test}: passed, four tests.
  \item \texttt{flutter build web}: passed; existing dependency Wasm dry-run
  warnings from \texttt{desktop\_drop}, \texttt{speech\_to\_text}, and
  \texttt{image} were reported.
  \item Touched Dart line counts: 104, 48, 201, 117, and 221 lines.
  \item Extracted-component usage grep found the new matters components
  referenced from the matters view and list composition.
\end{itemize}

\section*{Remaining Files at or Above 430 Lines}
\begin{itemize}
  \item 471: \texttt{lib/app/components/ai\_assistant/chat\_area.dart}
  \item 460: \texttt{lib/app/components/portal/calendar\_view.dart}
  \item 440: \texttt{lib/app/components/blog/blog\_feed.dart}
\end{itemize}

\end{document}
